\section{Motivation}
\label{sec:navhal_motivation}

Embedded flight control systems require precise timing, efficient execution, and direct interaction with hardware peripherals. The choice of hardware abstraction layer (HAL) plays a crucial role in achieving these requirements.

High-level frameworks such as the Arduino ecosystem prioritize ease of development. They provide simple and intuitive APIs that enable rapid prototyping. However, this abstraction hides low-level details, limiting control over hardware behavior and introducing execution overhead that is difficult to quantify. This lack of visibility and control makes such frameworks unsuitable for time-critical control systems.

In contrast, vendor-provided HALs expose detailed access to peripherals and enable efficient use of hardware resources. While this allows high-performance implementations, these frameworks are often complex, tightly coupled to specific microcontroller families, and difficult to scale or port across platforms.

These two approaches represent opposing design philosophies—one optimized for simplicity and the other for control and performance—leaving a gap in between that is not adequately addressed by existing solutions.

\begin{table}[H]
\centering
\begin{tabular}{|l|c|c|c|}
\hline
\textbf{Feature} & \textbf{Arduino Framework} & \textbf{Vendor HAL} & \textbf{NavHAL} \\
\hline
Ease of Use & High & Low & High \\
\hline
Performance & Low & High & High \\
\hline
Hardware Control & Limited & Extensive & Controlled \\
\hline
Portability & Moderate & Low & High \\
\hline
Determinism & Low & Moderate & High \\
\hline
Abstraction Overhead & High & Moderate & Low \\
\hline
Learning Curve & Low & High & Moderate \\
\hline
\end{tabular}
\caption{Comparison of abstraction approaches highlighting differences in usability, control, and performance.}
\label{tab:hal_comparison}
\end{table}

NavHAL is designed to address this gap by providing a structured interface that maintains low-level control while simplifying interaction with hardware. It achieves this through compile-time configuration, consistent APIs, and direct mapping to underlying peripherals.

The design enables predictable execution, reduced overhead, and improved portability without introducing additional runtime complexity. At the same time, it remains flexible enough to operate independently in bare-metal systems or alongside higher-level execution frameworks such as VAIOS.